\chapter*{ওপেন লজিক প্রকল্প সম্পর্কে}
\addcontentsline{toc}{chapter}{ওপেন লজিক প্রকল্প সম্পর্কে}


\textit{ওপেন লজিক পাঠ} হল আনুষ্ঠানিক অধিযুক্তিবিদ্যা ও আনুষ্ঠানিক
পদ্ধতির একটি উন্মুক্ত উৎসের, সহযোগিতায় রচিত পাঠ্যবই। এটি মধ্যবর্তী
স্তর থেকে শুরু হয় (অর্থাৎ, আনুষ্ঠানিক যুক্তিবিদ্যার প্রারম্ভিক
পাঠক্রমের পরে)। গণিতবিশেষজ্ঞ নন এমন পাঠকদের, বিশেষত দর্শন ও
কম্পিউটার বিজ্ঞানের শিক্ষার্থীদের, জন্য লেখা হলেও এটি যুক্তিগতভাবে
কঠোর।

বর্তমানে অন্তর্ভুক্ত কিছু বিষয়ের আলোচনা এখনও অসম্পূর্ণ হতে পারে,
এবং বহু বিভাগের উল্লেখযোগ্য সংশোধন প্রয়োজন। ভবিষ্যতে আরও বিষয়
অন্তর্ভুক্ত করে পাঠটি সম্প্রসারণের পরিকল্পনা আছে। পরিভাষাকোষ,
অতিরিক্ত পাঠের তালিকা, ঐতিহাসিক টীকা, ছবি, আরও ভাল ব্যাখ্যা,
দর্শন, কম্পিউটার বিজ্ঞান ও গণিতে ফলগুলির প্রাসঙ্গিকতা ব্যাখ্যাকারী
বিভাগ, এবং আরও সমস্যা ও উদাহরণ যোগ করারও পরিকল্পনা আছে।
ভুল পেলে বা পরামর্শ থাকলে
\href{https://github.com/OpenLogicProject/OpenLogic/wiki/Contributing}{প্রকল্পের দলকে জানাও}।

প্রকল্পটি উন্মুক্ত উৎসের আদর্শে পরিচালিত হয়। পাঠটি বিনামূল্যে
পাওয়া যায়; সঙ্গে ক্রিয়েটিভ কমন্স অ্যাট্রিবিউশন লাইসেন্সে
LaTeX উৎসও দেওয়া হয়। যথাযথ স্বীকৃতি দিলে এই লাইসেন্স যে কাউকে
কাজটি ডাউনলোড, ব্যবহার, পরিবর্তন, পুনর্বিন্যাস, রূপান্তর ও পুনর্বিতরণ
করার অধিকার দেয়। আরও তথ্যের জন্য ওপেন লজিক প্রকল্পের ওয়েবসাইট
\href{http://openlogicproject.org/}{openlogicproject.org} দেখো।
